\documentclass[10pt]{article}

\usepackage[letterpaper,margin=0.72in]{geometry}
\usepackage{booktabs}
\usepackage{array}
\usepackage{xcolor}
\usepackage{parskip}

\definecolor{midgray}{gray}{0.40}

% Compact lists (no enumitem available)
\setlength{\leftmargini}{1.5em}
\newcommand{\heading}[1]{\vspace{4pt}\noindent{\large\bfseries #1}\par\vspace{1pt}}
\newcommand{\notetext}[1]{{\footnotesize\color{midgray} #1\par}}

\setlength{\parskip}{3pt}
\pagestyle{empty}

\begin{document}

{\LARGE\bfseries Funding Proposal: GPU Compute Node for the Dartmouth Research Cluster\par}
\vspace{3pt}
{\color{black}\hrule height 1pt}
\vspace{2pt}

\heading{Summary}
We propose a one-time investment of \$185{,}000 to purchase a high-end GPU compute node and add
it to the existing Dartmouth research cluster (Slurm-scheduled, InfiniBand-connected). The node
provides four NVIDIA H200 accelerators with 564~GB of pooled high-bandwidth memory. It addresses
a gap in large-memory accelerated computing that currently sends AI and machine-learning research
to metered cloud GPUs. As a shared cluster resource, it would serve multiple labs while supporting
model training, fine-tuning, and high-throughput inference, and would help keep Dartmouth
competitive in AI research.

\heading{Justification}
\begin{itemize}
  \setlength{\itemsep}{2pt}\setlength{\parsep}{0pt}\setlength{\topsep}{2pt}
  \item \textbf{Capability gap.} Training, fine-tuning, and serving current open models (roughly
    70B to 400B parameters) need 100~GB or more of GPU memory and NVLink-coupled accelerators. The
    cluster has no large-VRAM GPUs today, so this work either waits in queue or runs on costly
    on-demand cloud instances.
  \item \textbf{Cost relative to cloud.} A comparable four-GPU H200 cloud instance rents for about
    \$14 to \$16 per hour, or roughly \$90{,}000 to \$120{,}000 per year at sustained research use.
    The node reaches break-even in about 18 to 24 months and then provides two to three further
    years of compute at only power and maintenance cost, on campus, with no data leaving the
    institution.
  \item \textbf{Competitiveness in AI research.} Large-memory GPUs are now standard infrastructure
    for competitive AI and machine-learning research. On-site capacity lets faculty and students
    train and fine-tune models and run inference at scale, and supports the externally funded AI
    projects that increasingly assume access to this class of hardware. It strengthens Dartmouth's
    position for faculty recruiting, grant competitiveness, and collaborations.
  \item \textbf{Shared resource.} The node joins the existing Slurm scheduler, so idle cycles are
    available to other groups across campus.
\end{itemize}

\heading{Proposed System and Budget}
\begin{center}
\renewcommand{\arraystretch}{1.08}
\begin{tabular}{@{}p{1.45in} p{4.05in} >{\raggedleft\arraybackslash}p{0.80in}@{}}
\toprule
\textbf{Component} & \textbf{Specification} & \textbf{Cost (USD)} \\
\midrule
GPU accelerators & 4 $\times$ NVIDIA H200 SXM, 141~GB HBM3e (HGX baseboard, NVLink/NVSwitch); 564~GB pooled & \$128{,}000 \\
Server platform & Dual AMD EPYC (192 cores), HGX carrier, redundant 3~kW PSUs & \$20{,}000 \\
System memory & 1.5~TB DDR5 ECC & \$10{,}000 \\
Local storage & About 30~TB NVMe (RAID1 OS plus high-speed scratch) & \$8{,}000 \\
Cluster networking & NVIDIA ConnectX-7 NDR 400~Gb InfiniBand HCA plus cabling & \$6{,}000 \\
Integration & Rack, rails, PDU, Slurm integration, burn-in testing & \$3{,}000 \\
Warranty and support & 3-year 24/7 on-site, next-business-day parts & \$10{,}000 \\
\midrule
\textbf{Total} & & \textbf{\$185{,}000} \\
\bottomrule
\end{tabular}
\end{center}

\notetext{Cost figures are planning estimates from our June 2026 conversation with Dartmouth
Research Computing (ITC) and are subject to final vendor quotes. Power and cooling are assumed
available in the existing cluster facility.}
\vspace{1pt}
\notetext{Alternative configuration: for workloads dominated by many-user inference rather than
large-model training, the same budget buys 8 $\times$ NVIDIA L40S (48~GB) for higher aggregate
throughput. The H200 build above is recommended as the more training-capable option.}

\heading{Expected Outcomes}
\begin{itemize}
  \setlength{\itemsep}{2pt}\setlength{\parsep}{0pt}\setlength{\topsep}{2pt}
  \item Continuous on-campus training and inference for AI and machine-learning research, without
    queueing for scarce accelerators.
  \item Reduced recurring cloud GPU spend once the node is in sustained use.
  \item Capacity to train and fine-tune models on institutional data, kept on campus.
  \item A shared accelerator resource that strengthens Dartmouth's competitiveness for AI research,
    recruiting, and grants.
\end{itemize}

\end{document}
